本文针对六足机器人 MPC 中一个隐蔽但影响极大的失效模式进行了系统修复：优化器可能在任务空间看似合理的情况下，收敛到物理上错误的 tarsus 分支。我们证明，该问题的根因并不主要来自步态、代价权重或 runtime guard 逻辑，而在于原有 SQP 路径没有在求解器内部充分执行分支可行性约束。为解决这一问题，本文引入硬约束 \emph{TarsusBranchConstraint}，并采用内点法（IPM）求解重构后的最优控制问题。

实验结果表明，所提控制器在报告的闭环测试中消除了 branch-guard 触发，恢复了稳定的前向推进，提高了跟踪精度，并始终保持在实时控制所需的计算预算之内。除短程 A/B 对照外，控制器还在中长时域 endurance 评测中保持了受控的偏航漂移、有限的机身高度波动以及亚 5~ms 的平均求解时间。比性能改进更重要的是架构层面的转变：物理分支有效性现在被纳入优化问题本身，而不再依赖优化后的 runtime 补救。这意味着控制系统从“事后止损”转向了“事前消除危险解”的设计范式。

本文的更广泛意义在于，高维多足系统中的求解器设计不能仅以目标函数收敛和算时为评价标准，还必须以其是否保持机构层面的物理语义为标准。在这类系统中，局部极小值可能在数值上完全无害，却在物理上不可接受。本文结果表明，将 IPM 与显式的机制约束相结合，是避免此类陷阱、构建更强安全解释能力的多足机器人 MPC 的有效方向。

总而言之，本文表明，求解器可见的分支一致性约束并不是一种可有可无的数值技巧，而是构建可靠六足机器人在线 MPC 的前提条件。通过将这类机制约束与 IPM 后端结合，并明确记录保证实时复现所需的工程条件，我们给出了一种同时具备科学可解释性与工程可部署性的控制设计路径。
